The ground outside King’s Landing did not want to be dug. Brienne had learned that much within the first hour of setting men to it, watching pick after pick ring uselessly off frozen clay packed hard as brick by a winter that showed no sign yet of loosening its grip, and had understood, standing there in the grey morning cold, that Sansa Stark’s plan — sound as it was, and she did believe it sound — was going to cost them considerably more in broken tools and blistered hands than anyone at that war council fire had properly reckoned for.

“Fires first,” she told the work-captains, because she had learned this too, watching sieges as a girl before she’d ever been permitted to fight in one. “Burn the topsoil soft before you set a blade to it. We haven’t the days to spend breaking picks on ground that isn’t ready for them.”

They worked in shifts through what remained of that day and the three that followed it, a thin, exhausted line of Northern levies and Unsullied and whatever smallfolk still had the strength left in their arms, digging a staggered ring of trenches and stake-lines across the open ground between the kingsroad’s final approach and King’s Landing’s walls — not a proper siege line, because they hadn’t the numbers left for anything a maester would recognize by that name in a book, but a killing ground, shaped the way Jaime Lannister had shown Theon Greyjoy at Winterfell in what felt now like some other lifetime: narrow, forced approaches, spears braced behind stakes, dragonglass held back for whatever got past the spears.

Jaime walked the lines with her most mornings, quiet in a way she had learned, across a year of knowing him properly rather than merely admiring him from a distance, meant he was thinking about his sister rather than the fortifications in front of him, and Brienne found she did not press him on it, because a man three days from his own sister’s shut gates was entitled to whatever silence he needed to carry that particular weight properly.

“You’ve an eye for this,” he told her, on the second morning, watching her redirect a line of stakes a few degrees east to close a gap she’d noticed no one else had. “Renly never taught you siege-craft. I’d wager my sword hand on it — the man barely troubled himself to learn it properly himself.”

“I taught myself. Watching, mostly. Reading, when there was time for it, which there rarely was.” She did not look up from the stake-line, because she had also learned, this past year, that compliments from Jaime Lannister still unsettled her in a way she preferred not examining too closely in front of him. “A woman who wants to be taken seriously with a sword learns early that the sword is only ever half of it. The other half is knowing where to stand so the sword matters at all.”

“That’s a harder lesson than swordwork. Most knights never trouble to learn it.” He was quiet a moment, watching a work-party wrestle a stubborn stake into ground that fought them for every inch. “Tormund would have liked this. The scale of it. He always did have a taste for building things men could die usefully behind, rather than dying prettily in front of.”

She let that sit a moment, feeling the old grief move through her the way it still did, months on — not sharp anymore, not the raw thing it had been at Winterfell’s southern gate, but present regardless, settled into her the way a scar settles into skin that has otherwise healed around it. “He’d have complained about the cold the whole while doing it.”

“Loudly,” Jaime agreed, and something in his face, for a moment, came as close to a real smile as she’d seen from him since they’d come within sight of these walls.

\sceneBreak

Rowan Glover worked the eastern stake-line with his own household men, competent and quiet and, Brienne thought, watching him direct a crew of exhausted diggers with more patience than she’d expected from a lord who’d buried his father on this same retreat, plainly grieving still in whatever private way a man allowed himself when there was no time left over for grieving properly. She had no cause to think of him as anything other than what he appeared to be — a young lord doing hard, necessary work without complaint, the kind of man a besieged army needed considerably more of than it presently had — and did not know, watching him, what Tyrion had begun, alone at the edge of camp some nights past, to suspect of someone wearing exactly that description.

“Lord Glover,” she said, coming to inspect his section on the third morning, “your line holds better than most. Where did you learn stake-work?”

“My father taught me, Ser Brienne. Deepwood Motte’s held against ironborn raiders longer than most Northern keeps have needed to bother learning how.” He did not quite meet her eyes saying it, though she read the flinch, if she read it at all, as simple grief rather than anything stranger — a young man still unaccustomed to being addressed by his father’s absent title. “I’d rather be useful with a spade than useless with a sword, given how little either’s likely to matter once the real fighting starts.”

She found nothing in the exchange worth remarking on further, and moved on down the line, and would not think of it again until considerably later, when the shape of what she’d actually been looking at came clear to her in hindsight the way most treacheries did — obvious only once a person already knew to look for it.

\sceneBreak

Grey Worm drove the Unsullied lines with a discipline that unsettled even Brienne to watch some mornings — men digging past the point where their hands had stopped properly obeying them, because their commander had not once, in three days, permitted himself to call a halt before the work itself was finished, and none of them, watching him, seemed inclined to be the first to ask for one either. She found him alone once, past midnight on the second night, standing over a half-finished stake-line with a torch held steady in a hand that did not shake despite hours of labor behind it, and understood, watching his face in the torchlight, that he had found in this work the same thing she had found in her own — a narrowness that kept the larger grief at bay, hour by hour, trench by trench.

“You’ll work them past what they’ve left to give,” she said, not unkindly, “if you don’t call the shift.”

“They have more to give than you know, Ser Brienne. I have more to give than you know.” He did not look at her, still measuring the stake-line’s angle against the tree line beyond it. “Missandei asked me once whether I was afraid of anything anymore. I told her fear was cut from me long before I had the words to understand what was being taken. I think now that was only half true. I am not afraid of the dead, or the dying. I find I am considerably more afraid of finishing this war with nothing left in my hands to build with them.” He set the torch into the ground at last, satisfied, and turned to face her properly. “So I dig, Ser Brienne. It is the one thing this war has not yet found a way to take from me.”

She had no answer worth offering to that, and gave him none, and left him to his lines.

The food ran shorter with each of those three days than any war council wanted said plainly at its own fire, though everyone at it understood the arithmetic regardless — a column that had marched into starvation across six hundred miles of retreat could not simply stop marching and expect the stopping to feed it, and Davos’s reports each evening grew leaner in the telling the same way the grain stores grew leaner in fact, until the smallfolk labor gangs were working trench-lines on half rations and the Dothraki’s last horses, the few that remained, had begun disappearing from their picket lines into cookpots no one troubled to ask too many questions about. Brienne ate what she was given and did not ask for more, because she had learned, watching men dig twelve hours on a bowl of thin barley broth, that complaining about her own share bought nothing back for anyone standing in that trench beside her.

The walls above them stayed shut and watchful the whole three days, Golden Company banners moving along the battlements in slow, deliberate patrols that Brienne understood, watching them, were meant less for defense than for demonstration — a queen behind stone making certain the starving army at her gates understood precisely how comfortable that stone kept her, and precisely how little comfort it meant to extend past its own walls. Once, on the second day, a handful of city archers loosed a scattering of arrows down into the work parties, not enough to kill more than two men outright, and Brienne understood that too as demonstration rather than earnest attack — Cersei reminding them, in the smallest possible increments, exactly how many ways she could hurt them without opening a single gate to do it.

Daenerys did not fly Drogon over the walls again after the first day, once Tyrion’s warning about the scorpions’ spacing had made its way properly around the war council, and Brienne found herself missing the dragon’s presence overhead more than she’d expected to — not for what he could have done against the walls, since even Tyrion agreed one dragon could not break a city this size alone, but for the simple, animal comfort of something enormous and alive still standing between this exhausted camp and whatever came for it next.

\sceneBreak

The scout came in at dusk on the fourth day, a Northern rider Brienne recognized from Sansa’s own household guard, his horse blown and lathered in a way that told her the news before he’d said a word of it, the particular exhausted urgency of a man who had ridden considerably harder than the message strictly required, because some part of him understood that arriving a quarter hour later might have meant arriving too late to matter.

“They’ve closed the distance,” he said, breathless, to the assembled war council Brienne had been summoned to along with the rest. “Faster than the last count reckoned. We took it for a fortnight, three days past. It isn’t. Six days, maybe five, if the weather holds against us instead of them.”

“How,” Jon said, flat, and Brienne heard in the single word the same question every face around that fire was already asking.

“Melisandre’s fire, near as the scouts could tell it,” the rider said. “The trench-line she lit at Winterfell’s edge — it’s held them up this whole retreat, slower going round the burn scar than the direct road would’ve let them travel. It’s spent now. Burned through, or they’ve found a way round it at last, and either way it isn’t costing them the days it used to.” He swallowed, and Brienne watched him steady himself before finishing it. “Whatever time her dying bought us, my lords, my ladies — I think we’ve spent the last of it.”

No one spoke for a long moment after that, and Brienne, watching the fire’s light move across faces gone still with the particular arithmetic of people recalculating exactly how much time remained to finish work that had already felt, three days ago, like considerably more work than the time allotted for it could properly hold.

“Then we work through the nights as well as the days,” Sansa said, finally, quiet, the flat calculating stillness Tyrion had remarked on more than once these past weeks settling fully over her face. “Every trench unfinished by dawn five days from now is a trench that never mattered at all. Ser Brienne — I want the eastern line’s gaps closed before anything else. That’s where I’d come through, were I commanding what’s coming for us.”

“It’ll be done,” Brienne said, and found, saying it, that some old, familiar steadiness settled into her the way it always had whenever duty narrowed itself down to a single, achievable thing rather than the vast, unmanageable grief of the war entire — a wall to build, a gap to close, a line to hold, and however small a comfort that narrowing offered against everything still waiting to arrive out of the dark behind them, she found she was grateful for it regardless, the same plain, unglamorous gratitude a woman learns for whatever solid ground remains beneath her feet, when so much else has already proven it could not be trusted to hold.

She walked the eastern line herself that same night, torch in hand, measuring gaps by eye and by memory both, and did not let herself look north, toward the dark kingsroad where six days, or five, were presently closing the last of whatever distance remained between this exhausted camp and the thing that had already unmade a wall seven hundred years old and a castle that had stood since the Age of Heroes, because there was work still to finish before dawn, and a besieged army, she had learned, survived considerably longer on work finished than on fear properly reckoned with.